% Options for packages loaded elsewhere
\PassOptionsToPackage{unicode}{hyperref}
\PassOptionsToPackage{hyphens}{url}
\documentclass[
]{article}
\usepackage{xcolor}
\usepackage{amsmath,amssymb}
\setcounter{secnumdepth}{5}
\usepackage{iftex}
\ifPDFTeX
  \usepackage[T1]{fontenc}
  \usepackage[utf8]{inputenc}
  \usepackage{textcomp} % provide euro and other symbols
\else % if luatex or xetex
  \usepackage{unicode-math} % this also loads fontspec
  \defaultfontfeatures{Scale=MatchLowercase}
  \defaultfontfeatures[\rmfamily]{Ligatures=TeX,Scale=1}
\fi
% \usepackage{lmodern}
\ifPDFTeX\else
  % xetex/luatex font selection
\fi
% Use upquote if available, for straight quotes in verbatim environments
\IfFileExists{upquote.sty}{\usepackage{upquote}}{}
\IfFileExists{microtype.sty}{% use microtype if available
  \usepackage[]{microtype}
  \UseMicrotypeSet[protrusion]{basicmath} % disable protrusion for tt fonts
}{}
\makeatletter
\@ifundefined{KOMAClassName}{% if non-KOMA class
  \IfFileExists{parskip.sty}{%
    \usepackage{parskip}
  }{% else
    \setlength{\parindent}{0pt}
    \setlength{\parskip}{6pt plus 2pt minus 1pt}}
}{% if KOMA class
  \KOMAoptions{parskip=half}}
\makeatother
\usepackage{color}
\usepackage{fancyvrb}
\newcommand{\VerbBar}{|}
\newcommand{\VERB}{\Verb[commandchars=\\\{\}]}
\DefineVerbatimEnvironment{Highlighting}{Verbatim}{commandchars=\\\{\}}
% Add ',fontsize=\small' for more characters per line
\newenvironment{Shaded}{}{}
\newcommand{\AlertTok}[1]{\textcolor[rgb]{1.00,0.00,0.00}{\textbf{#1}}}
\newcommand{\AnnotationTok}[1]{\textcolor[rgb]{0.38,0.63,0.69}{\textbf{\textit{#1}}}}
\newcommand{\AttributeTok}[1]{\textcolor[rgb]{0.49,0.56,0.16}{#1}}
\newcommand{\BaseNTok}[1]{\textcolor[rgb]{0.25,0.63,0.44}{#1}}
\newcommand{\BuiltInTok}[1]{\textcolor[rgb]{0.00,0.50,0.00}{#1}}
\newcommand{\CharTok}[1]{\textcolor[rgb]{0.25,0.44,0.63}{#1}}
\newcommand{\CommentTok}[1]{\textcolor[rgb]{0.38,0.63,0.69}{\textit{#1}}}
\newcommand{\CommentVarTok}[1]{\textcolor[rgb]{0.38,0.63,0.69}{\textbf{\textit{#1}}}}
\newcommand{\ConstantTok}[1]{\textcolor[rgb]{0.53,0.00,0.00}{#1}}
\newcommand{\ControlFlowTok}[1]{\textcolor[rgb]{0.00,0.44,0.13}{\textbf{#1}}}
\newcommand{\DataTypeTok}[1]{\textcolor[rgb]{0.56,0.13,0.00}{#1}}
\newcommand{\DecValTok}[1]{\textcolor[rgb]{0.25,0.63,0.44}{#1}}
\newcommand{\DocumentationTok}[1]{\textcolor[rgb]{0.73,0.13,0.13}{\textit{#1}}}
\newcommand{\ErrorTok}[1]{\textcolor[rgb]{1.00,0.00,0.00}{\textbf{#1}}}
\newcommand{\ExtensionTok}[1]{#1}
\newcommand{\FloatTok}[1]{\textcolor[rgb]{0.25,0.63,0.44}{#1}}
\newcommand{\FunctionTok}[1]{\textcolor[rgb]{0.02,0.16,0.49}{#1}}
\newcommand{\ImportTok}[1]{\textcolor[rgb]{0.00,0.50,0.00}{\textbf{#1}}}
\newcommand{\InformationTok}[1]{\textcolor[rgb]{0.38,0.63,0.69}{\textbf{\textit{#1}}}}
\newcommand{\KeywordTok}[1]{\textcolor[rgb]{0.00,0.44,0.13}{\textbf{#1}}}
\newcommand{\NormalTok}[1]{#1}
\newcommand{\OperatorTok}[1]{\textcolor[rgb]{0.40,0.40,0.40}{#1}}
\newcommand{\OtherTok}[1]{\textcolor[rgb]{0.00,0.44,0.13}{#1}}
\newcommand{\PreprocessorTok}[1]{\textcolor[rgb]{0.74,0.48,0.00}{#1}}
\newcommand{\RegionMarkerTok}[1]{#1}
\newcommand{\SpecialCharTok}[1]{\textcolor[rgb]{0.25,0.44,0.63}{#1}}
\newcommand{\SpecialStringTok}[1]{\textcolor[rgb]{0.73,0.40,0.53}{#1}}
\newcommand{\StringTok}[1]{\textcolor[rgb]{0.25,0.44,0.63}{#1}}
\newcommand{\VariableTok}[1]{\textcolor[rgb]{0.10,0.09,0.49}{#1}}
\newcommand{\VerbatimStringTok}[1]{\textcolor[rgb]{0.25,0.44,0.63}{#1}}
\newcommand{\WarningTok}[1]{\textcolor[rgb]{0.38,0.63,0.69}{\textbf{\textit{#1}}}}
\usepackage{longtable,booktabs,array}
\newcounter{none} % for unnumbered tables
\usepackage{calc} % for calculating minipage widths
% Correct order of tables after \paragraph or \subparagraph
\usepackage{etoolbox}
\makeatletter
\patchcmd\longtable{\par}{\if@noskipsec\mbox{}\fi\par}{}{}
\makeatother
% Allow footnotes in longtable head/foot
\IfFileExists{footnotehyper.sty}{\usepackage{footnotehyper}}{\usepackage{footnote}}
\makesavenoteenv{longtable}
\usepackage{graphicx}
\makeatletter
\newsavebox\pandoc@box
\newcommand*\pandocbounded[1]{% scales image to fit in text height/width
  \sbox\pandoc@box{#1}%
  \Gscale@div\@tempa{\textheight}{\dimexpr\ht\pandoc@box+\dp\pandoc@box\relax}%
  \Gscale@div\@tempb{\linewidth}{\wd\pandoc@box}%
  \ifdim\@tempb\p@<\@tempa\p@\let\@tempa\@tempb\fi% select the smaller of both
  \ifdim\@tempa\p@<\p@\scalebox{\@tempa}{\usebox\pandoc@box}%
  \else\usebox{\pandoc@box}%
  \fi%
}
% Set default figure placement to htbp
\def\fps@figure{htbp}
\makeatother
\setlength{\emergencystretch}{3em} % prevent overfull lines
\providecommand{\tightlist}{%
  \setlength{\itemsep}{0pt}\setlength{\parskip}{0pt}}
\usepackage[]{natbib}
\bibliographystyle{plainnat}
\makeatletter
\@ifpackageloaded{subcaption}{}{\usepackage{subcaption}}
\@ifpackageloaded{caption}{}{\usepackage{caption}}
\captionsetup[subfigure]{margin=0.5em}
\AtBeginDocument{%
\renewcommand*\figurename{Figure}
\renewcommand*\tablename{Table}
}
\AtBeginDocument{%
\renewcommand*\listfigurename{List of Figures}
\renewcommand*\listtablename{List of Tables}
}
\newcounter{pandoccrossref@subfigures@footnote@counter}
\newenvironment{pandoccrossrefsubfigures}{%
\setcounter{pandoccrossref@subfigures@footnote@counter}{0}
\begin{figure}\centering%
\gdef\global@pandoccrossref@subfigures@footnotes{}%
\DeclareRobustCommand{\footnote}[1]{\footnotemark%
\stepcounter{pandoccrossref@subfigures@footnote@counter}%
\ifx\global@pandoccrossref@subfigures@footnotes\empty%
\gdef\global@pandoccrossref@subfigures@footnotes{{##1}}%
\else%
\g@addto@macro\global@pandoccrossref@subfigures@footnotes{, {##1}}%
\fi}}%
{\end{figure}%
\addtocounter{footnote}{-\value{pandoccrossref@subfigures@footnote@counter}}
\@for\f:=\global@pandoccrossref@subfigures@footnotes\do{\stepcounter{footnote}\footnotetext{\f}}%
\gdef\global@pandoccrossref@subfigures@footnotes{}}
\@ifpackageloaded{float}{}{\usepackage{float}}
\floatstyle{ruled}
\@ifundefined{c@chapter}{\newfloat{codelisting}{h}{lop}}{\newfloat{codelisting}{h}{lop}[chapter]}
\floatname{codelisting}{Listing}
\newcommand*\listoflistings{\listof{codelisting}{List of Listings}}
\makeatother
\usepackage{bookmark}
\IfFileExists{xurl.sty}{\usepackage{xurl}}{} % add URL line breaks if available
\urlstyle{same}
% fallback for those not using the hyperref driver hyperxmp:
\makeatletter
\@ifundefined{xmpquote}{\newcommand{\xmpquote}[1]{#1}}{}
\makeatother
\hypersetup{
  pdftitle={Registered Report Template: Preregistration{,} Deviations{,} and Claim Boundaries},
  pdfauthor={Daniel Ari Friedman},
  colorlinks=true,linkcolor=red,urlcolor=red,citecolor=red,anchorcolor=red,filecolor=red,
  pdfcreator={LaTeX via pandoc}}

\title{Registered Report Template: Preregistration, Deviations, and
Claim Boundaries}
\author{Daniel Ari Friedman}
\date{}

\IfFileExists{seqsplit.sty}{\usepackage{seqsplit}}{\newcommand{\seqsplit}[1]{#1}}
\protected\def\breakseq#1{\seqsplit{#1}}
\protected\def\breaktt#1{\begingroup\ttfamily\seqsplit{#1}\endgroup}
\title{Registered Report Template: Preregistration, Deviations, and Claim Boundaries}
\author{Daniel Ari Friedman\\\footnotesize{Active Inference Institute}\\\footnotesize{\texttt{daniel@activeinference.institute}}\\\footnotesize{\href{https://orcid.org/0000-0001-6232-9096}{ORCID: 0000-0001-6232-9096}} \\ \footnotesize{DOI: \href{https://doi.org/10.5281/zenodo.21298892}{10.5281/zenodo.21298892}}}
\date{2026-07-10}

\usepackage[margin=0.75in]{geometry}
\usepackage{graphicx}
\setlength{\parindent}{0pt}
\setlength{\parskip}{0.45em}

\begin{document}
\begin{titlepage}
\centering
\vspace*{2cm}
{\LARGE\bfseries Registered Report Template: Preregistration, Deviations, and Claim Boundaries\par}
\vspace{0.75em}
{\large A reproducible exemplar for registered analysis plans\par}
\vfill
\makeatletter
{\@author\par}
\vspace{1em}
{\@date\par}
\makeatother
\vfill
\end{titlepage}
\thispagestyle{empty}
\tableofcontents
\newpage

\section{Abstract}\label{sec:abstract}

This document is a \textbf{template}, not an empirical study. It
demonstrates the registered-report workflow end to end: locking a
preregistration, validating its completeness, executing the registered
analysis plan against deterministic demonstration data, and reporting
confirmatory and exploratory claims through an explicit deviation
ledger. Every quantity reported here is produced by the tested code in
\breaktt{src/registered\_report/} and regenerated by
\breaktt{scripts/generate\_figures.py}; none is hand-entered or
illustrative.

The demonstration binds a single confirmatory hypothesis (H1) to one
registered outcome (\texttt{primary\_score}) analysed by a two-sided
label-permutation test. Run on a seeded two-group dataset
(\breaktt{seed\ =\ 20260709}, \texttt{n\ =\ 24} per group), the
registered test yields an observed mean difference of \texttt{1.003}
with a two-sided permutation p-value of \texttt{0.0005} (\texttt{0} of
\texttt{2000} shuffles at least as extreme), significant at the
preregistered \texttt{alpha\ =\ 0.05}. A deliberately introduced
secondary endpoint and an alternative model are carried only as
\textbf{documented deviations}, keeping the confirmatory claim boundary
intact. The purpose is to give forks a working, auditable skeleton in
which planned analyses, frozen hypotheses, deviations, and post-run
claims are separable and machine-checkable.

\newpage

\section{Introduction}\label{sec:introduction}

A registered report separates \emph{what a study commits to} from
\emph{what the data later show}. Hypotheses, outcomes, exclusion rules,
and the analysis plan are locked and time-stamped before results are
observed; only then is the analysis run and interpreted. This ordering
is the structural defence against the degrees-of-freedom problems that
inflate false-positive rates in flexible analyses
\citep{simmons2011false} and motivates the broader preregistration and
research-transparency movement
\citep{nosek2018preregistration, munafo2017manifesto, chambers2013registered}.

This template operationalises that discipline as tested code rather than
prose convention. The \breaktt{registered\_report} package freezes a
registration payload under a content hash, validates that the required
preregistration sections are present and well formed, and --- after
execution --- compares the analysis that was actually run against the
plan that was locked. Any departure must appear in a \textbf{deviation
ledger} with a rationale before it can be reported, and confirmatory
claims are kept mechanically distinct from exploratory ones.

To make the workflow concrete and reproducible, the template ships a
\textbf{deterministic demonstration study}. It contains no real
empirical data: a seeded generator synthesises a two-group dataset, and
the registered label-permutation test is executed against it. Because
both the data and the permutation schedule are seeded, every figure and
every number in this document is byte-stable and is regenerated on
demand from \breaktt{scripts/generate\_figures.py}. Readers forking this
template should replace the demonstration fixture in
\breaktt{data/example\_registration.json} and the synthetic data
generator with their own registration and real data collection, keeping
the same lock-validate-execute-compare structure.

The remainder of this report states the preregistered hypotheses
(Section \ref{sec:hypotheses}), describes the frozen methods and
analysis plan (Section \ref{sec:methods}), reports the confirmatory
outcome bound to the executed analysis (Section \ref{sec:results}),
records deviations (Section \ref{sec:deviations}), and discusses scope
and boundaries (Section \ref{sec:discussion}).

\newpage

\section{Preregistered hypotheses}\label{sec:hypotheses}

The registration in \breaktt{data/example\_registration.json} declares
its hypotheses, outcomes, exclusion rules, and analysis plan before any
result is observed. It is frozen under a SHA-256 content hash (prefix
\texttt{96a34a11d132}); the validator in
\breaktt{registered\_report.protocol.validate\_registration} confirms the
payload carries every required section with no structural findings.

\subsection{Confirmatory hypothesis}\label{confirmatory-hypothesis}

\begin{itemize}
\tightlist
\item
  \textbf{H1} --- \emph{The proposed condition improves the primary
  score.} This is the sole confirmatory hypothesis. It maps to the
  registered outcome \texttt{primary\_score}, measured as a mean score
  difference and analysed with a two-sided label-permutation test at
  \texttt{alpha\ =\ 0.05}.
\end{itemize}

Each registered hypothesis is bound to exactly one registered outcome
and one registered analysis. fig.~\ref{fig:hypothesis_map} renders that
binding directly from the frozen registration via
\breaktt{hypothesis\_outcome\_map}, so the figure cannot silently drift
from the locked plan.

\begin{figure}
\centering
\pandocbounded{\includegraphics[keepaspectratio,alt={Registered hypothesis-to-outcome-to-analysis mapping for the demonstration registration. Each row is generated by hypothesis\_outcome\_map(registration) in src/registered\_report/demo\_study.py: hypothesis H1 links to the primary\_score outcome, the mean score difference measure, and the two-sided permutation test analysis. Only registered hypotheses appear; exploratory endpoints are excluded by construction.}]{../figures/hypothesis_map.png}}
\caption{Registered hypothesis-to-outcome-to-analysis mapping for the
demonstration registration. Each row is generated by
\protect\breaktt{hypothesis\_outcome\_map(registration)} in
\protect\breaktt{src/registered\_report/demo\_study.py}: hypothesis H1
links to the \protect\texttt{primary\_score} outcome, the
\protect\breaktt{mean\ score\ difference} measure, and the
\protect\breaktt{two-sided\ permutation\ test} analysis. Only registered
hypotheses appear; exploratory endpoints are excluded by
construction.}\label{fig:hypothesis_map}
\end{figure}

\subsection{Registered outcome and exclusion
rule}\label{registered-outcome-and-exclusion-rule}

\begin{itemize}
\tightlist
\item
  \textbf{Primary outcome} \texttt{primary\_score} --- mean score
  difference between the proposed and reference conditions, evaluated by
  a two-sided permutation test.
\item
  \textbf{Exclusion rule} --- records missing the primary outcome are
  excluded \emph{before} group assignment is inspected, so exclusion
  decisions cannot depend on the observed effect.
\end{itemize}

No secondary outcome is preregistered. Any additional endpoint
introduced during analysis is therefore exploratory and must be recorded
in the deviation ledger (Section \ref{sec:deviations}) rather than
presented as a confirmatory finding.

\newpage

\section{Methods and analysis plan}\label{sec:methods}

The method is the registered-report workflow itself, implemented as
tested code. It has four locked stages that complete before any result
is observed, followed by two interpretation stages that run only after
the deviation ledger is built. fig.~\ref{fig:analysis_dag} shows this
ordering and the lock boundary it protects.

\begin{figure}
\centering
\pandocbounded{\includegraphics[keepaspectratio,alt={Registered-report analysis workflow, rendered from analysis\_plan\_stages() in src/registered\_report/demo\_study.py. The four blue stages --- freeze registration, validate completeness, collect/simulate data, and execute the registered analysis --- lock before results are revealed (dashed boundary). The two purple stages --- compare to plan and build the deviation ledger, then interpret confirmatory claims --- run only afterwards. The locks\_before\_results flag on each stage is what enforces the boundary in code.}]{../figures/analysis_dag.png}}
\caption{Registered-report analysis workflow, rendered from
\protect\breaktt{analysis\_plan\_stages()} in
\protect\breaktt{src/registered\_report/demo\_study.py}. The four blue
stages --- freeze registration, validate completeness, collect/simulate
data, and execute the registered analysis --- lock before results are
revealed (dashed boundary). The two purple stages --- compare to plan
and build the deviation ledger, then interpret confirmatory claims ---
run only afterwards. The \protect\breaktt{locks\_before\_results} flag on
each stage is what enforces the boundary in
code.}\label{fig:analysis_dag}
\end{figure}

\subsection{Frozen registration}\label{frozen-registration}

\breaktt{freeze\_registration} deep-copies the registration, strips any
pre-existing hash, and stamps a canonical SHA-256 hash over the
sorted-key JSON encoding. \breaktt{validate\_registration} then checks
that \texttt{title}, \texttt{version}, \texttt{hypotheses},
\texttt{outcomes}, \texttt{exclusion\_rules}, and
\texttt{analysis\_plan} are present; that each hypothesis has a unique
\texttt{id} and a \texttt{claim}; that each outcome names a
\texttt{measure} and \texttt{analysis}; that
\breaktt{analysis\_plan.primary\_model} and \texttt{seed} are set; and
that registered-report stage and ethics-review metadata are supplied.
For the demonstration registration this validation returns no findings.

\subsection{Registered analysis plan}\label{registered-analysis-plan}

The analysis plan is fixed before execution:

\begin{itemize}
\tightlist
\item
  \textbf{Primary model} --- a two-sided label-permutation test on the
  group mean difference.
\item
  \textbf{Significance threshold} --- \texttt{alpha\ =\ 0.05}.
\item
  \textbf{Seed} --- \texttt{20260709}, used for both data synthesis and
  the permutation schedule so the result is reproducible.
\item
  \textbf{Permutations} --- \texttt{2000} label shuffles; the two-sided
  p-value uses the add-one correction
  \breaktt{(k\ +\ 1)\ /\ (permutations\ +\ 1)} and is therefore strictly
  positive.
\end{itemize}

\subsection{Deterministic demonstration
data}\label{deterministic-demonstration-data}

Because this is a template, the ``data collection'' stage is a seeded
synthetic generator (\breaktt{generate\_demo\_dataset}), \textbf{not} a
real study. It draws \texttt{n\ =\ 24} control observations from
\texttt{Normal(0,\ 1)} and \texttt{n\ =\ 24} treatment observations from
\texttt{Normal(0.8,\ 1)} using a single seeded generator, so the entire
dataset is a pure function of the plan seed.
\breaktt{run\_registered\_analysis} reads the seed and \texttt{alpha}
back out of the frozen plan and executes the registered permutation test
against this dataset, guaranteeing the executed analysis honours the
locked plan. The executed summary is written to
\breaktt{output/data/demo\_analysis.json} and is the single source of
truth for every number in Section \ref{sec:results}.

Forks replace this generator with genuine data collection while keeping
the same \breaktt{run\_registered\_analysis} contract, so the
confirmatory analysis remains bound to the frozen plan rather than to
post-hoc choices.

\newpage

\section{Results}\label{sec:results}

Results are interpreted only after the registered analysis has run and
the deviation ledger (Section \ref{sec:deviations}) has classified any
departures from the plan. Every value below is read from
\breaktt{output/data/demo\_analysis.json}, which is regenerated by
\breaktt{scripts/generate\_figures.py} from the tested
\breaktt{run\_registered\_analysis} function.

\subsection{\texorpdfstring{Confirmatory outcome (H1,
\texttt{primary\_score})}{Confirmatory outcome (H1, primary_score)}}\label{confirmatory-outcome-h1-primary_score}

On the deterministic demonstration data, the registered two-sided
permutation test returns:

{\def\LTcaptype{none} % do not increment counter
\begin{longtable}[]{@{}ll@{}}
\toprule\noalign{}
Quantity & Value \\
\midrule\noalign{}
\endhead
\bottomrule\noalign{}
\endlastfoot
Control mean & \texttt{-0.178991} \\
Treatment mean & \texttt{0.823948} \\
Observed mean difference & \texttt{1.002940} \\
Permutations & \texttt{2000} \\
Shuffles with absolute difference \texorpdfstring{\textgreater{}=}{>=} observed & \texttt{0} \\
Two-sided permutation p-value & \texttt{0.0005} \\
Preregistered \texttt{alpha} & \texttt{0.05} \\
Confirmatory decision & significant \\
\end{longtable}
}

The observed difference of \texttt{1.003} lies far in the right tail of
the permutation null distribution: none of the \texttt{2000} label
shuffles produced a mean difference as extreme, so the add-one-corrected
p-value is \texttt{0.0005}, below the preregistered
\texttt{alpha\ =\ 0.05}. H1 is therefore \textbf{supported} as a
confirmatory result. fig.~\ref{fig:permutation_result} shows the
observed statistic against the null it was tested against.

\begin{figure}
\centering
\pandocbounded{\includegraphics[keepaspectratio,alt={Registered permutation test on the deterministic demonstration data, rendered by plot\_permutation\_result from permutation\_null\_distribution and run\_permutation\_test. The histogram is the label-permutation null distribution of the treatment-minus-control mean difference (2000 seeded shuffles); the solid red line marks the observed difference of 1.003 and the dotted line its mirror image. The annotation reports n = 24 per group, 0 shuffles at least as extreme, a two-sided p-value of 0.0005, and alpha = 0.05. The observed statistic sits outside the entire null, which is why the primary confirmatory claim holds.}]{../figures/permutation_result.png}}
\caption{Registered permutation test on the deterministic demonstration
data, rendered by \protect\breaktt{plot\_permutation\_result} from
\protect\breaktt{permutation\_null\_distribution} and
\protect\breaktt{run\_permutation\_test}. The histogram is the
label-permutation null distribution of the treatment-minus-control mean
difference (\protect\texttt{2000} seeded shuffles); the solid red line
marks the observed difference of \protect\texttt{1.003} and the dotted
line its mirror image. The annotation reports \protect\texttt{n\ =\ 24}
per group, \protect\texttt{0} shuffles at least as extreme, a two-sided
p-value of \protect\texttt{0.0005}, and \protect\texttt{alpha\ =\ 0.05}.
The observed statistic sits outside the entire null, which is why the
primary confirmatory claim holds.}\label{fig:permutation_result}
\end{figure}

\subsection{Claim boundary}\label{claim-boundary}

The registration preregisters exactly one confirmatory outcome. The
review packet built by \breaktt{build\_review\_packet} records
\texttt{primary\_score} as the sole confirmatory outcome and returns an
integrity score of \texttt{1.0} for a plan-faithful execution. Any other
endpoint examined during analysis is exploratory and is reported under
Section \ref{sec:deviations}, never here. This separation is enforced in
code, not merely asserted in prose.

\newpage

\section{Deviation register}\label{sec:deviations}

A registered report must state, explicitly, every way the executed
analysis departed from the frozen plan. This template treats that
register as a first-class artifact: \breaktt{build\_deviation\_ledger}
produces one row per executed outcome and model, classifying each as
\texttt{ok} (registered), \texttt{warning} (an unregistered element
carried with a documented rationale), or \texttt{error} (an unregistered
element with no rationale). \breaktt{compare\_analysis\_to\_registration}
downgrades a documented model change from \texttt{error} to
\texttt{warning} and refuses to accept any deviation whose
\texttt{rationale} is empty.

\subsection{Demonstration deviations}\label{demonstration-deviations}

To exercise the machinery, the demonstration executes a deliberately
plan-divergent analysis: it adds a \texttt{secondary\_score} endpoint
and swaps the registered \breaktt{permutation\_test} for a
\texttt{linear\_model}, each accompanied by a rationale. The resulting
ledger is:

{\def\LTcaptype{none} % do not increment counter
\begin{longtable}[]{@{}llllll@{}}
\toprule\noalign{}
Order & Kind & Target & Registered & Documented & Severity \\
\midrule\noalign{}
\endhead
\bottomrule\noalign{}
\endlastfoot
0 & outcome & \texttt{primary\_score} & yes & --- & \texttt{ok} \\
1 & outcome & \texttt{secondary\_score} & no & yes & \texttt{warning} \\
2 & model & \texttt{linear\_model} & no & yes & \texttt{warning} \\
\end{longtable}
}

Because both departures are documented, the review packet remains valid:
it reports \texttt{primary\_score} as the only confirmatory outcome,
\texttt{secondary\_score} as exploratory, and an integrity score of
\texttt{0.9} (a single \texttt{0.10} warning penalty against a perfect
plan). Had either deviation lacked a rationale, the same code would have
raised a blocking \breaktt{deviation\_without\_rationale} or
\breaktt{unregistered\_outcome} error, and the packet would be invalid.
fig.~\ref{fig:deviation_timeline} plots this register along the
registered-to-executed axis.

\begin{figure}
\centering
\pandocbounded{\includegraphics[keepaspectratio,alt={Preregistration-versus-executed deviation ledger timeline, rendered by plot\_deviation\_timeline from build\_deviation\_ledger and deviation\_timeline. Each marker is a ledger row coloured by severity: the registered primary\_score outcome is green (ok), while the added secondary\_score endpoint and the linear\_model substitution are amber (warning, documented deviations). No marker is red because every departure carries a rationale. An undocumented departure would appear in red and invalidate the review packet.}]{../figures/deviation_timeline.png}}
\caption{Preregistration-versus-executed deviation ledger timeline,
rendered by \protect\breaktt{plot\_deviation\_timeline} from
\protect\breaktt{build\_deviation\_ledger} and
\protect\breaktt{deviation\_timeline}. Each marker is a ledger row
coloured by severity: the registered \protect\texttt{primary\_score}
outcome is green (\protect\texttt{ok}), while the added
\protect\texttt{secondary\_score} endpoint and the
\protect\texttt{linear\_model} substitution are amber
(\protect\texttt{warning}, documented deviations). No marker is red
because every departure carries a rationale. An undocumented departure
would appear in red and invalidate the review
packet.}\label{fig:deviation_timeline}
\end{figure}

\subsection{Reading the register}\label{reading-the-register}

The register is the mechanism that keeps Section \ref{sec:results}
honest: the confirmatory claim is restricted to the registered
\texttt{primary\_score} outcome, and the exploratory
\texttt{secondary\_score} endpoint is visible, labelled, and excluded
from confirmatory interpretation. Forks should extend this ledger ---
never edit the frozen registration --- whenever the executed analysis
must diverge from plan.

\newpage

\section{Discussion}\label{sec:discussion}

\subsection{What the demonstration does and does not
show}\label{what-the-demonstration-does-and-does-not-show}

The demonstration shows that the registered analysis, run against seeded
data with a genuine effect (\texttt{Normal(0.8,\ 1)} treatment versus
\texttt{Normal(0,\ 1)} control), recovers a significant confirmatory
result: an observed mean difference of \texttt{1.003} with a two-sided
permutation p-value of \texttt{0.0005}. It does \textbf{not} show
anything about any real-world phenomenon. The data are synthetic and the
effect is injected by construction; the number exists to prove that the
locked plan, when executed, produces an auditable statistic --- not to
support a scientific claim. This is the honest reading of a template:
the value is a property of the code path, not evidence about the world.

\subsection{Why the boundaries are enforced in
code}\label{why-the-boundaries-are-enforced-in-code}

Prose conventions for ``confirmatory versus exploratory'' are easy to
blur once results are in view. Here the boundary is a function contract:
\breaktt{compare\_analysis\_to\_registration} refuses unregistered
outcomes without a documented deviation,
\breaktt{build\_deviation\_ledger} assigns a severity to every executed
element, and \breaktt{build\_review\_packet} partitions confirmatory from
exploratory outcomes deterministically. The content hash on the frozen
registration makes silent edits to the locked plan detectable
(\texttt{hash\_mismatch}). Together these turn registered-report
discipline into checks that fail loudly rather than guidance that can be
quietly ignored.

\subsection{Limitations and scope}\label{limitations-and-scope}

\begin{itemize}
\tightlist
\item
  The demonstration uses a single hypothesis and a single confirmatory
  outcome; real registrations will declare several, and the same helpers
  scale to them.
\item
  A permutation test is used because it is exact, deterministic, and
  dependency-free; forks may register any primary model by naming it in
  the plan.
\item
  Ethics-review and registered-report-stage metadata in the fixture are
  synthetic (\texttt{status:\ exempt},
  \breaktt{authority:\ synthetic-fixture-review}) and must be replaced
  with real values before any submission.
\end{itemize}

\subsection{Using this template}\label{using-this-template}

Fork with \breaktt{scripts/audit/copy\_exemplar.py}, replace
\breaktt{data/example\_registration.json} with your own preregistration,
swap the synthetic generator in
\breaktt{src/registered\_report/demo\_study.py} for real data collection,
record stage and ethics metadata, and rerun the tests and figure
generation. Keep confirmatory claims tied to registered outcomes, record
every departure in the deviation ledger, and never let post-run result
prose rewrite the preregistered intent.

\newpage

\section{References}\label{sec:references}

The bibliography lives in
\href{references.bib}{\breaktt{manuscript/references.bib}} and is read by
Pandoc during PDF rendering. Every \texttt{{[}@key{]}} citation in the
manuscript is resolved against that file, so the reference list below is
generated from the cited entries rather than maintained by hand.

To validate that \texttt{references.bib} is syntactically clean and has
the required fields per entry type, run from the repository root:

\begin{Shaded}
\begin{Highlighting}[]
\ExtensionTok{uv}\NormalTok{ run python }\AttributeTok{{-}m}\NormalTok{ infrastructure.reference.citation.cli validate }\DataTypeTok{\textbackslash{}}
\NormalTok{    projects/templates/template\_registered\_report/manuscript/references.bib }\AttributeTok{{-}{-}strict}
\end{Highlighting}
\end{Shaded}




\clearpage

\bibliography{references}
\end{document}
